\section{관련 연구}
\label{sec:related}

\subsection{IMU 자세 추정}

Mahony 등~\cite{mahony2008}은 SO(3) 위에서의 비선형 상보 필터를 제안하여,
가속도계 및 지자기 센서 오차로 구동되는 비례-적분 피드백 법칙을 이용해 자이로스코프 적분을 보정하였다.
Madgwick 등~\cite{madgwick2011}은 센서 측정값과 예측 방향 사이의 오차를 최소화하는
경사 하강 기반 필터를 제안하여, 명시적인 노이즈 모델 없이 단일 튜닝 파라미터만으로
Mahony 필터에 필적하는 정확도를 달성하였다.
EKF에 기반한 확률적 접근법은 현재 추정값 주위에서 쿼터니언 운동학을 선형화하고
완전한 공분산을 전파한다~\cite{bar-shalom2001}.
ESKF~\cite{sola2017}는 4성분 쿼터니언 섭동 대신 최소한의 3성분 오차 상태 회전 벡터를 유지함으로써
EKF를 개선하여 특이점을 방지하고 선형화 오차를 줄인다.
이러한 필터 계열의 비교 평가는 보행자 항법~\cite{woodman2007}과 로봇 플랫폼에서 수행되었으나,
명시적인 지터 분석을 포함한 아두이노급 UAV 하드웨어에서의 체계적인 비교는 상대적으로 드물다.

\subsection{UAV 비행 제어 및 상태 추정}

멀티로터 자세 안정화를 위한 계단식 PID 제어는 잘 확립되어 있다~\cite{bouabdallah2004}.
PX4 자동 조종 장치~\cite{px4docs}는 2단계 계단식 제어기를 구현하는데,
외부 자세 루프가 내부 각속도 루프에서 소비되는 동체 각속도 설정값을 생성하며,
이를 커뮤니티에서는 일반적으로 PPID라고 부른다.
ArduPilot~\cite{ardupilotdocs}은 피드포워드 및 노치 필터링 기능이 추가된 구조적으로 유사한 아키텍처를 따른다.
두 프레임워크 모두 PPID 제어기를 IMU, 기압계, GPS, 광학 흐름 데이터를 융합하는
다중 주파수 EKF(ArduPilot의 EKF2/EKF3, PX4의 EKF2)와 결합한다~\cite{mahony2008}.
MuJoCo~\cite{todorov2012}와 Gazebo 같은 물리 엔진을 이용한 이러한 파이프라인의 시뮬레이션 기반 검증은
자동 조종 장치 개발의 표준 관행이 되었으며,
Google DeepMind의 MuJoCo Menagerie~\cite{deepmind2022}는
이 목적에 적합한 고충실도 항공기 모델을 제공한다.
